% Pista alg-24j-q2 · Àlgebra
% Procedència: PAU Catalunya 2024 · Convocatòria de juny · Sèrie 1 · Qüestió 2
\documentclass[11pt,a4paper]{article}
\usepackage{pista}
\pistainfo{Catalunya 2024}{Convocatòria de juny}{Sèrie 1}

\begin{document}

\section*{\textcolor{blauPista}{Qüestió 2}}

\noindent Considereu el sistema d'equacions següent:
\[
\left.\begin{array}{r}
4x + 2y - z = 4 \\
x - y + kz = 3 \\
3x + 3y = 1
\end{array}\right\},
\]
on $k$ és un paràmetre real.

\subsection*{\textit{a)} \normalfont Discutiu el sistema per als diferents valors del paràmetre $k$, i resoleu-lo per a $k = 0$. \hfill\textnormal{[1~punt]}}

\pista{Fixa't que et demanen dues coses ben diferents, i en una sola frase. Et demanen "discutiu" i després, "resoleu". Per tant, per discutir quin tipus de sistema serà, escriu la matriu de coeficients $A$ i la matriu ampliada $A^{'}$. Calcula $\det(A)$ i troba els valors del paràmetre $k$ que l'anul·len. Per a cadascun d'aquests valors, compara els rangs de $A$ i $A^{'}$. Per la segona part de la frase, has de resoldre per a $k = 0$. Com que aquest valor no serà cap dels que has trobar a l'apartat anterior que anul·laven el determinant, vol dir que el sistema serà compatible determinat i, per tant, trobaràs una única solució numèrica.}

\subsection*{\textit{b)} \normalfont Resoleu el sistema per a $k = -1$. \hfill\textnormal{[0,75~punts]}}

\pista{Aquest valor és, precisament, el que anul·la el determinant de $A$. Analitza amb cura els rangs d'$A$ i $A^{'}$. El cas més complicat acostuma a ser quan el sistema sigui compatible indeterminat. En aquest cas, va bé introduir una variable lliure (per exemple $x = \lambda$) i expressar la resta de variables en funció d'ella. És a dir, $y$ i $z$ quedaran en funció del paràmetre $\lambda$.}

\subsection*{\textit{c)} \normalfont Per a $k = -1$, modifiqueu la tercera equació de manera que el sistema esdevingui incompatible. Justifiqueu la resposta. \hfill\textnormal{[0,75~punts]}}

\pista{Has d'aconsguir que els rangs de $A$ i $A^{'}$ no coincideixin, perquè d'aquesta manera el sistema serà incompatible.}

\end{document}
